\documentclass[UTF8,fontset=fandol,11pt]{ctexart}
\usepackage[a4paper,margin=24mm]{geometry}
\usepackage{amsmath,amssymb,amsthm,graphicx,tabularx,booktabs,hyperref}
\newtheorem{definition}{定义}[section]
\newtheorem{theorem}[definition]{定理}
\newtheorem{example}[definition]{例题}
\newtheorem{exercise}[definition]{练习}
\hypersetup{unicode=true,colorlinks=true,linkcolor=blue}
\title{最小二乘法讲义}
\date{}
\begin{document}
\maketitle
\section{误差如何度量：从残差到常数拟合}\label{sec:S01}
% UNIT S01.U01
对同一个量重复测量时，读数往往并不完全相同。如果要用一个数代表这些读数，应当如何选择？最小二乘法先规定误差的度量，再寻找使这种误差最小的模型参数。本节从常数拟合入手，只使用有限求和、一元求导和配方法。\par
% UNIT S01.U02
\begin{definition}
设 $n\geq 1$ 为整数，$y_i\in\mathbb{R}$ 是第 $i$ 次观测的响应值，$i=1,\ldots,n$。若记录中还有自变量，则用 $x_i\in\mathbb{R}$ 表示第 $i$ 次观测的输入。给定模型及其参数后，模型对第 $i$ 次观测给出的数值记为 $\hat{y}_i$，称为拟合值。

约定残差为观测值减去拟合值：
\[
r_i=y_i-\hat{y}_i,\qquad i=1,\ldots,n.
\]
因此，$r_i>0$ 表示模型低估了该观测值，$r_i<0$ 表示模型高估了该观测值。残差平方和定义为
\[
Q=\sum_{i=1}^{n}r_i^2=\sum_{i=1}^{n}(y_i-\hat{y}_i)^2.
\]
最小二乘准则是：在选定模型的允许参数范围内，寻找使 $Q$ 取得最小值的参数。这里最小化的是所选模型对已有数据的残差平方和，而不是保证每个残差都分别最小。
\end{definition}
% UNIT S01.U03
为什么不直接用残差之和度量拟合误差？正负残差可能互相抵消。例如，两个残差为 $1$ 与 $-1$ 时，残差之和为零，但两个拟合值都没有命中观测值。平方后各项非负，不能发生这种抵消；并且
\[
\sum_{i=1}^{n}r_i^2=0\quad\Longleftrightarrow\quad r_i=0\ \text{对所有 }i=1,\ldots,n.
\]
平方也会使较大的残差更显著地影响目标函数。如果响应值有计量单位，残差与响应值同单位，而残差平方和具有该单位的平方。比较不同候选模型的平方和时，本节始终使用同一组数据和相同单位。\par
% UNIT S01.U04
\begin{definition}
设 $n\geq1$ 为整数，观测值为 $y_1,\ldots,y_n\in\mathbb{R}$。常数模型用同一个待定实数 $c$ 拟合所有观测值，即
\[
\hat{y}_i=c,\qquad i=1,\ldots,n,\qquad c\in\mathbb{R}.
\]
相应的目标函数是
\[
Q(c)=\sum_{i=1}^{n}(y_i-c)^2,\qquad c\in\mathbb{R}.
\]
这些观测值的样本算术平均数记为
\[
\bar{y}=\frac{1}{n}\sum_{i=1}^{n}y_i.
\]
在优化过程中，观测值以及由它们算出的 $\bar{y}$ 都是固定的，只有参数 $c$ 发生变化。
\end{definition}
% UNIT S01.U05
\begin{theorem}
设 $n\geq1$ 为整数，$y_1,\ldots,y_n\in\mathbb{R}$，并令 $\bar{y}=\frac{1}{n}\sum_{i=1}^{n}y_i$。则常数模型的残差平方和满足：对每个 $c\in\mathbb{R}$，
\[
Q(c)=\sum_{i=1}^{n}(y_i-\bar{y})^2+n(c-\bar{y})^2.
\]
因此，$Q(c)$ 在且仅在 $c=\bar{y}$ 时取得全局最小值，最小值为
\[
Q(\bar{y})=\sum_{i=1}^{n}(y_i-\bar{y})^2.
\]
\end{theorem}
% UNIT S01.U06
\begin{proof}
先用一元求导寻找候选参数。对任意 $c\in\mathbb{R}$，有限项逐项求导得
\[
Q'(c)=2\sum_{i=1}^{n}(c-y_i)=2n(c-\bar{y}).
\]
所以 $Q'(c)=0$ 当且仅当 $c=\bar{y}$。由于参数范围是整个实数轴，可微函数若在某点取得最小值，该点必须满足导数为零。不过，求出驻点本身还不是全局最小性的证明。

下面直接比较所有参数对应的函数值。由平均数的定义，
\[
\sum_{i=1}^{n}(y_i-\bar{y})=\sum_{i=1}^{n}y_i-n\bar{y}=0.
\]
于是，对任意 $c\in\mathbb{R}$，平方展开得到
\[
\begin{aligned}
Q(c)&=\sum_{i=1}^{n}\bigl((y_i-\bar{y})+(\bar{y}-c)\bigr)^2\\
&=\sum_{i=1}^{n}(y_i-\bar{y})^2
 +2(\bar{y}-c)\sum_{i=1}^{n}(y_i-\bar{y})
 +n(\bar{y}-c)^2\\
&=Q(\bar{y})+n(c-\bar{y})^2.
\end{aligned}
\]
因为 $n>0$，最后一项对所有实数 $c$ 都非负，且仅在 $c=\bar{y}$ 时为零。这既证明了最小值存在，也证明了全局最小点唯一。
\end{proof}
% UNIT S01.U07
\begin{example}
某个保持不变的量经过三次测量，读数依次为 $2,3,7$，单位相同。用常数模型拟合这些读数：写出目标函数，求最优常数及最小残差平方和，并与候选常数 $c=3$ 比较。
\end{example}
% UNIT S01.U08
\par\noindent\textbf{解：}对 $c\in\mathbb{R}$，目标函数为
\[
\begin{aligned}
Q(c)&=(2-c)^2+(3-c)^2+(7-c)^2\\
&=3c^2-24c+62\\
&=3(c-4)^2+14.
\end{aligned}
\]
三个读数的平均数是 $\bar{y}=(2+3+7)/3=4$，所以唯一最优常数为 $c=4$。这也可从 $Q'(c)=6c-24$ 和上述配方看出。

将两个候选常数的残差列在一起：
\[
\begin{array}{c|rrr|r}
\text{候选常数}&r_1&r_2&r_3&Q(c)\\ \hline
c=3&-1&0&4&17\\
c=4&-2&-1&3&14
\end{array}
\]
因此最小平方和为 $Q(4)=4+1+9=14$，而 $Q(3)=1+0+16=17$。从 $c=3$ 改为 $c=4$ 后，前两个残差的绝对值反而增大，但总平方和减小了 $3$。这说明最小二乘法优化的是整体指标，并不要求每次测量的拟合误差都同时减小。\par
% UNIT S01.U09
常数拟合给出了一个可重复使用的思路：先明确参数与目标函数，再求导寻找候选点，最后用平方和分解证明全局最优性。恒等式 $Q(c)-Q(\bar{y})=n(c-\bar{y})^2$ 还定量说明了偏离最优参数的代价。

最优常数模型的残差之和为 $\sum_{i=1}^{n}(y_i-\bar{y})=0$，但残差平方和不一定为零；它为零当且仅当所有观测值相同。第\ref{sec:S02}节将把这里的常数拟合结论用于直线拟合，仍然只需要一元求导与配方。\par
% UNIT S01.U10
\begin{exercise}
\label{mathagent:exercise:79e91357ee05b8b9aa4c657d553157f1ef411624d282c3f2337dd0e01453ce2d}
对同单位的四次测量值 $1,2,4,5$ 作常数拟合。写出 $Q(c)$ 的平方和形式，求最优常数、各个残差及最小平方和；再计算 $Q(2)$，并用平方和分解解释它与最小值之差。本题答案统一置于讲义末尾。
\end{exercise}
% UNIT S01.U11
\begin{exercise}
\label{mathagent:exercise:9e7dbf7ef17df1506f236bbf1ad6ff8c14ff13065df89e12b796d18475228be7}
设 $n\geq1$ 为整数，$y_1,\ldots,y_n\in\mathbb{R}$，常数模型的残差为 $r_i=y_i-c$，其中 $c\in\mathbb{R}$。判断并说明理由：

（1）若 $\sum_{i=1}^{n}r_i=0$，则 $c$ 是该常数模型的唯一最优参数。

（2）若 $\sum_{i=1}^{n}r_i=0$，则每一个残差都等于零。

（3）当 $n=1$ 时，常数模型仍有唯一最优参数。

本题答案统一置于讲义末尾。
\end{exercise}
\section*{练习答案}
% UNIT S01.U10.answer
\par\noindent\textbf{\hyperref[mathagent:exercise:79e91357ee05b8b9aa4c657d553157f1ef411624d282c3f2337dd0e01453ce2d]{练习 \ref*{mathagent:exercise:79e91357ee05b8b9aa4c657d553157f1ef411624d282c3f2337dd0e01453ce2d} 的解答：}}目标函数为
\[
Q(c)=(1-c)^2+(2-c)^2+(4-c)^2+(5-c)^2,\qquad c\in\mathbb{R}.
\]
四次测量值的平均数为
\[
\bar{y}=\frac{1+2+4+5}{4}=3.
\]
由常数拟合的平方和分解，
\[
\begin{aligned}
Q(c)&=(1-3)^2+(2-3)^2+(4-3)^2+(5-3)^2+4(c-3)^2\\
&=10+4(c-3)^2.
\end{aligned}
\]
因此唯一最优常数为 $c=3$，对应的残差依次为
\[
r_1=-2,\qquad r_2=-1,\qquad r_3=1,\qquad r_4=2,
\]
最小平方和为 $Q(3)=4+1+1+4=10$。

当 $c=2$ 时，
\[
Q(2)=(-1)^2+0^2+2^2+3^2=14.
\]
平方和分解给出
\[
Q(2)-Q(3)=4(2-3)^2=4.
\]
也就是说，候选常数 $2$ 偏离平均数 $3$ 一个单位，使残差平方和比最小值增加了 $4$。\par
% UNIT S01.U11.answer
\par\noindent\textbf{\hyperref[mathagent:exercise:9e7dbf7ef17df1506f236bbf1ad6ff8c14ff13065df89e12b796d18475228be7]{练习 \ref*{mathagent:exercise:9e7dbf7ef17df1506f236bbf1ad6ff8c14ff13065df89e12b796d18475228be7} 的解答：}}（1）正确。由残差定义与平均数定义，
\[
\sum_{i=1}^{n}r_i=\sum_{i=1}^{n}y_i-nc=n(\bar{y}-c).
\]
因为 $n>0$，残差之和为零必有 $c=\bar{y}$。由常数拟合定理，$\bar{y}$ 正是该模型的唯一最优参数。

（2）错误。残差之和为零可以由正负残差抵消产生。例如取 $n=2$、$y_1=1$、$y_2=3$ 及 $c=2$，则
\[
r_1=-1,\qquad r_2=1,\qquad r_1+r_2=0,
\]
但两个残差均不为零。因此题设不能推出每个残差都为零。

（3）正确。当 $n=1$ 时，
\[
Q(c)=(y_1-c)^2\geq0.
\]
等号当且仅当 $c=y_1$ 成立，所以唯一最优参数为 $c=y_1=\bar{y}$，最小平方和为零。\par
\end{document}
